\begin{table}[H]
\centering
\caption{Peer/Rank Diagnostic and Classroom Shifts: Liberia}
\label{tab:suffstat_liberia}
\begin{threeparttable}
\begin{tabular}[t]{lc}
\toprule
\multicolumn{1}{c}{ } & \multicolumn{1}{c}{Liberia} \\
\cmidrule(l{3pt}r{3pt}){2-2}
\multicolumn{1}{c}{ } & \multicolumn{1}{c}{(1)} \\
\cmidrule(l{3pt}r{3pt}){2-2}
 & Estimate \\
\midrule
\multicolumn{2}{l}{\textit{Panel A: Approximate Borusyak--Hull peer/rank coefficient}} \\
$\hat{\zeta}_{Lib}$ &  -0.048 (  0.108) \\
Observations & 3153 \\
\addlinespace
\multicolumn{2}{l}{\textit{Panel B: Treatment-induced classroom shifts}} \\
$\Delta$ absolute EB-to-target mismatch (C $-$ T) &   0.348 \\
$\Delta$ peer composition (T $-$ C) &   0.041 \\
$\Delta$ class size (T $-$ C) &    11.4 \\
$\Delta$ within class baseline SD (C $-$ T) &   0.058 \\
\addlinespace
Overall ITT &  -0.211 (  0.135) \\
\bottomrule
\end{tabular}
\begin{tablenotes}[para,flushleft]
\footnotesize
\item \textit{Notes:} Panel A reports the approximate Borusyak--Hull control function peer/rank coefficient using decile $\times$ treatment $\times$ grade FE\@. Standard errors are clustered at the school grade group level. Panel B reports treatment-control differences in classroom characteristics. The near-zero peer/rank coefficient and small classroom composition shift provide little evidence that the measured peer/rank diagnostic accounts for the negative Liberia ITT pattern.
\end{tablenotes}
\end{threeparttable}
\end{table}
